% Generated draft: requires mathematical and source review.
\documentclass[11pt]{article}
\usepackage[margin=25mm]{geometry}
\usepackage{fontspec}
\setmainfont{Latin Modern Roman}
\usepackage{amsmath,amssymb,mathtools,unicode-math}
\setmathfont{Latin Modern Math}
\usepackage{enumitem,needspace,xurl,hyperref}
\setlength{\parindent}{0pt}
\setlength{\parskip}{5pt}
\setlength{\emergencystretch}{3em}
\newcommand{\N}{\mathbb{N}}
\newcommand{\Z}{\mathbb{Z}}
\newcommand{\Q}{\mathbb{Q}}
\newcommand{\R}{\mathbb{R}}
\newcommand{\C}{\mathbb{C}}
\newcommand{\F}{\mathbb{F}}
\newcommand{\A}{\mathbb{A}}
\newcommand{\PP}{\mathbb{P}}
\newcommand{\E}{\mathbb{E}}
\newcommand{\eps}{\varepsilon}
\newcommand{\dd}{\,\mathrm{d}}
\begin{document}
\section*{UnsolvedMath ID: 30003652}
\textbf{Source problem code:} catalogue entry
\subsection{Definitions and mathematical statement}
For a positive integer \(n\), let \(S^{n-1}=\{x\in\R^n:\|x\|_2=1\}\). The Euclidean kissing number is
\[
\tau_n=\max\bigl\{|C|:C\subseteq S^{n-1}\text{ is finite and }\langle x,y\rangle\leq \tfrac12\text{ for all distinct }x,y\in C\bigr\}.
\]
Equivalently, it is the largest number of unit balls with pairwise disjoint interiors that can simultaneously touch a fixed unit ball in \(\R^n\).

For terminology, a real symmetric \(m\times m\) matrix \(A\) is copositive if \(z^{\mathsf T}Az\geq0\) for every \(z\in\R^m\) with nonnegative coordinates. It is completely positive if \(A=BB^{\mathsf T}\) for some entrywise nonnegative real \(m\times r\) matrix \(B\), with \(r\) a positive integer. These definitions do not specify the approximation hierarchies intended by the source; in particular, the supplied literature assessment refers to copositive kernels without supplying their precise framework.

\textbf{Source question.} Determine the mathematical relationships between the copositive and completely positive approximation hierarchies referred to in the cited Oberwolfach report and other existing approximation hierarchies for bounding the kissing number.

\textbf{Unresolved formulation.} The original question identifies the first hierarchies only as those ``developed here.'' The supplied record does not give their cones, optimization problems, levels, or admissible kernels, and does not identify the other approximation hierarchies. Nor does it specify whether a relationship means inclusion of feasible regions, comparison of bounds at corresponding levels, equivalence after reindexing, or a relationship between limits. The definition of \(\tau_n\) above supplies the standard Euclidean meaning of kissing number; the intended dimensions and dimension quantifiers of the comparison remain unspecified. Consequently, the available evidence supports a comparison question, but does not determine a unique quantified conjecture or inequality.
\subsection{Short English statement}
The question asks how the copositive and completely positive hierarchies discussed in the source report relate to other methods for bounding the kissing number. The supplied record does not identify the hierarchies or the precise comparison sought, so a fully specified mathematical conjecture cannot be extracted. Its open status and reported source verification are supplied curation claims, not independent verification.
\subsection{Sources}
\begin{itemize}
\item Supplied local registry record: permanent UnsolvedMath ID \(30003652\), source problem code \texttt{catalogue entry}, title \emph{Copositive Hierarchies for the Kissing-Number Problem}, in the Oberwolfach Reports Open Problems collection. The record lists proposed year 2017 and supplies no individual proposer.
\item Supplied report citation: \emph{Copositivity and Complete Positivity} (2018), \emph{Oberwolfach Reports}, DOI 10.4171/owr/2017/52, pp.\ 3071--3120. \url{https://doi.org/10.4171/owr/2017/52}. These are the supplied report-wide bibliographic details; no problem-specific page is supplied.
\item Supplied literature reference: Kuryatnikova--Vera, approximating copositive kernels. \url{https://doi.org/10.1016/j.endm.2017.10.052}.
\item Supplied literature reference: Oberwolfach Report 52, open question. \url{https://doi.org/10.4171/owr/2017/52}.
\item Literature assessment dated exactly \texttt{2026-08-21}: the record labels the problem open and states that convergent copositive-kernel hierarchies for the kissing number have been developed, while the requested general relationship with other approximations remains unsettled. This is supplied dated curation, not an independently verified literature assessment; the record itself requests independent verification of time-sensitive status and cited claims.
\item The registry labels the statement \texttt{corrected\_verified} and reports that the official report was checked, with the source-dependent phrase ``developed here'' replaced by a direct comparison question. This verification account is supplied metadata; no independent source consultation is claimed in this draft.
\end{itemize}
\end{document}
